\section{Introduction}
\vspace{0.2in}
\hspace*{0.16in}
%TODO explain what we tested with steps For this chapter, we will present the tools we used and our results of segmenting and counting blood cells (red, white, platelets).

\section{Used Tools}
\subsection{TensorFlow}
In our work we used TensorFlow which is an open source library developed by Google Brain Team for Artificial Intelligence, it contains multiple pre-defined models and algorithms for Deep Learning and Machine Learning, tensorFlow can be used in multiple languages as python, C++, Java and JavaScript.\\
Tensorflow can be used to Create, Train, Deploy Models. And when we talk about complicated models we can use Keras.
Keras is a high level API for Neural Networks which helps with experiments on the models and extensibility.

\begin{figure}[H]
    \centering
      \vspace{-0.1in}
        \centerline{\includegraphics[width = 3in]{../images/tensorflow.png}}
        \caption{Tensorflow}
        \label{Tensorflow}
    \end{figure}

\subsection{Colab}
Google Colab or Colaboratory is a free Jupyter notebook environment running on Google's cloud servers for machine learning training and research. This platform allows the user to leverage backend hardware such as GPUs and TPUs and train Machine Learning and Deep Learning models directly in the cloud. Without the need to install anything on our computer at the anything on our computer except a browser.
but it has some disadvantages where we have a limit on the GPU usage
and non presistant storage.

\begin{figure}[H]
\centering
  \vspace{-0.1in}
    \centerline{\includegraphics[width = 2in]{../images/colab.png}}
    \caption{Google Colab}
    \label{Google Colab}
\end{figure}

\subsection{PaperSpace}
Paperspace is a high-performance cloud computing and ML development platform for building, training and deploying machine learning models. It has a complete jupyter notebook environement which has a persistent storage and 6 hours limit for each execution.

\begin{figure}[H]
    \centering
      \vspace{-0.1in}
        \centerline{\includegraphics[width = 3in]{../images/paperspace.png}}
        \caption{Paperspace}
        \label{Paperspace}
    \end{figure}

\section{Experiments and Results}

\subsection{Red Blood Cells}
\subsubsection{DO-U-Net}
%TODO we will have a table which has accuracy and metrics for example 10 images and their real and predicted count with the 3 methods
%example of table columns
% Image Results(Metrics IOU ...) Real-Count CHT CCL Watershed

\begin{table}[H]
\centering
\begin{tabular}{|  c | c | c | c | c | c | c | c |}
\hline
\textbf{Image} & \textbf{Real count} & \textbf{Watershed} & \textbf{CCL} & \textbf{CHT} & \textbf{Watershed\_acc} & \textbf{CCL\_acc} & \textbf{CHT\_acc} \\
\hline
Im037\_0 & 105 & 85 & 70 & 102 & 80.95 & 66.67 & 97.14 \\
Im045\_0 & 601 & 517 & 415 & 577 & 86.02 & 69.05 & 96.01 \\
Im053\_1 & 802 & 571 & 320 & 652 & 71.20 & 39.90 & 81.30 \\
Im001\_1 & 215 & 214 & 205 & 216 & 99.53 & 95.35 & 99.53 \\
Im004\_1 & 258 & 280 & 247 & 268 & 91.47 & 95.74 & 96.12 \\
Im015\_1 & 293 & 295 & 260 & 281 & 99.32 & 88.74 & 95.90 \\
Im022\_1 & 242 & 251 & 235 & 246 & 96.28 & 97.11 & 98.35 \\
Im050\_1 & 460 & 498 & 423 & 519 & 91.74 & 91.96 & 87.17 \\
Im069\_0 & 665 & 651 & 531 & 675 & 97.89 & 79.85 & 98.50 \\
Im079\_0 & 512 & 496 & 398 & 518 & 96.88 & 77.73 & 98.83 \\
Im095\_0 & 150 & 189 & 95 & 150 & 74.00 & 63.33 & 100.00 \\
Im099\_0 & 528 & 478 & 397 & 517 & 90.53 & 75.19 & 97.92 \\
Im108\_0 & 510 & 511 & 418 & 546 & 99.80 & 81.96 & 92.94 \\
\hline
\textbf{Total} &  &  &  &  & 90.43 & 78.66 & 95.36 \\
\hline
\end{tabular}
\caption{RBC Counting results using 3 algorithms}
\label{table:DOUNET-RBC}
\end{table}


\subsubsection{DO-SegNet}
%TODO we will do the same thing as DO-UNET and at the end we draw a resume table wich will compare between both/ or we can do both in the same table 

% \begin{table}[H]
\centering
\begin{tabular}{|c|c|c|c|c|c|c|c|}
\hline
 \textbf{Image} & \textbf{Real\_Count} & \textbf{CHT} & \textbf{CCL} & \textbf{EDT} & \textbf{CHT\_acc} & \textbf{CCL\_acc} & \textbf{EDT\_acc} \\ \hline
 Im001\_1 & 215 & 212 & 234 & 247 & 98.6 & 91.16 & 85.12 \\ 
 Im004\_1 & 258 & 255 & 264 & 285 & 98.84 & 97.67 & 89.53 \\ 
 Im015\_1 & 293 & 269 & 212 & 261 & 91.81 & 72.35 & 89.08 \\ 
 Im022\_1 & 242 & 232 & 279 & 268 & 95.87 & 84.71 & 89.26 \\ 
 Im037\_0 & 105 & 103 & 75 & 88 & 98.1 & 71.43 & 83.81 \\ 
 Im045\_0 & 601 & 576 & 409 & 480 & 95.84 & 68.05 & 79.87 \\ 
 Im050\_1 & 460 & 486 & 450 & 475 & 94.35 & 97.83 & 96.74 \\ 
 Im053\_1 & 802 & 614 & 228 & 448 & 76.56 & 28.43 & 55.86 \\ 
 Im069\_0 & 665 & 653 & 376 & 549 & 98.2 & 56.54 & 82.56 \\ 
 Im079\_0 & 512 & 502 & 407 & 430 & 98.05 & 79.49 & 83.98 \\ 
 Im095\_0 & 150 & 177 & 177 & 170 & 82.0 & 82.0 & 86.67 \\ 
 Im099\_0 & 528 & 516 & 437 & 465 & 97.73 & 82.77 & 88.07 \\ 
 Im108\_0 & 510 & 528 & 418 & 458 & 96.47 & 81.96 & 89.8 \\ \hline
 Total & -1 & -1 & -1 & -1 & 94.03 & 76.49 & 84.64 \\ 

\hline
\end{tabular}
\caption{Results of each algorithm}
\label{Results of each algorithm}
\end{table}


\subsection{White Blood Cells}
%TODO Same
\subsubsection{DO-U-Net}

\begin{table}[H]
\centering
\begin{tabular}{|  c | c | c | c | c | c | c | c |}
\hline
\textbf{Image} & \textbf{Real count} & \textbf{Watershed} & \textbf{CCL} & \textbf{CHT} & \textbf{Watershed\_acc} & \textbf{CCL\_acc} & \textbf{CHT\_acc} \\
\hline
Im037\_0 & 1 & 1 & 2 & 1 & 100.00 & 0.00 & 100.00 \\
Im045\_0 & 1 & 1 & 2 & 1 & 100.00 & 0.00 & 100.00 \\
Im053\_1 & 45 & 44 & 39 & 41 & 97.78 & 86.67 & 91.11 \\
Im001\_1 & 18 & 17 & 14 & 15 & 94.44 & 77.78 & 83.33 \\
Im004\_1 & 12 & 12 & 10 & 14 & 100.00 & 83.33 & 83.33 \\
Im015\_1 & 22 & 22 & 16 & 20 & 100.00 & 72.73 & 90.91 \\
Im022\_1 & 5 & 5 & 6 & 5 & 100.00 & 80.00 & 100.00 \\
Im050\_1 & 21 & 22 & 18 & 24 & 95.24 & 85.71 & 85.71 \\
Im069\_0 & 1 & 1 & 2 & 1 & 100.00 & 0.00 & 100.00 \\
Im079\_0 & 7 & 8 & 8 & 6 & 85.71 & 85.71 & 85.71 \\
Im095\_0 & 2 & 2 & 3 & 2 & 100.00 & 50.00 & 100.00 \\
Im099\_0 & 3 & 3 & 4 & 1 & 100.00 & 66.67 & 33.33 \\
Im108\_0 & 4 & 4 & 5 & 4 & 100.00 & 75.00 & 100.00 \\
\hline
\textbf{Total 13 images} &  &  &  &  & 97.94 & 58.74 & 88.7 \\
\textbf{Total 108 images} &  &  &  &  & 95.64 & 51.68 & 85.76 \\
\hline
\end{tabular}
\caption{WBC Counting results with 3 algorithms}
\label{table:UNet-WBC}
\end{table}

\subsubsection{DO-SegNet}

% \begin{table}[H]
    \centering
    \begin{tabular}{|c|c|c|c|c|c|c|c|}
    \hline
     \textbf{Image} & \textbf{Real\_Count} & \textbf{CHT} & \textbf{CCL} & \textbf{EDT} & \textbf{CHT\_acc} & \textbf{CCL\_acc} & \textbf{EDT\_acc} \\ \hline
     Im001\_1 & 18 & 13 & 13 & 16 & 72.22 & 72.22 & 88.89 \\ 
     Im004\_1 & 12 & 11 & 9 & 12 & 91.67 & 75.0 & 100.0 \\ 
     Im015\_1 & 22 & 18 & 15 & 21 & 81.82 & 68.18 & 95.45 \\ 
     Im022\_1 & 5 & 5 & 5 & 5 & 100.0 & 100.0 & 100.0 \\ 
     Im037\_0 & 1 & 0 & 1 & 1 & 0 & 100.0 & 100.0 \\ 
     Im045\_0 & 1 & 1 & 1 & 1 & 100.0 & 100.0 & 100.0 \\ 
     Im050\_1 & 21 & 29 & 20 & 21 & 61.9 & 95.24 & 100.0 \\ 
     Im053\_1 & 45 & 44 & 43 & 44 & 97.78 & 95.56 & 97.78 \\ 
     Im069\_0 & 1 & 1 & 1 & 1 & 100.0 & 100.0 & 100.0 \\ 
     Im079\_0 & 7 & 7 & 9 & 9 & 100.0 & 71.43 & 71.43 \\ 
     Im095\_0 & 2 & 2 & 2 & 2 & 100.0 & 100.0 & 100.0 \\ 
     Im099\_0 & 3 & 1 & 3 & 3 & 33.33 & 100.0 & 100.0 \\ 
     Im108\_0 & 4 & 4 & 8 & 4 & 100.0 & 0 & 100.0 \\ \hline
     Total & -1 & -1 & -1 & -1 & 79.9 & 82.89 & 96.43 \\ 
    
    \hline
    \end{tabular}
    \caption{Results of wbc segnet}
    \label{Results of wbc segnet}
    \end{table}
    

\subsection{Platlets}
%TODO Same
\subsubsection{DO-U-Net}
\subsubsection{DO-SegNet}

% \begin{table}[H]
\centering
\begin{tabular}{|c|c|c|c|c|c|c|c|}
\hline
 \textbf{Image} & \textbf{Real\_Count} & \textbf{CHT} & \textbf{CCL} & \textbf{EDT} & \textbf{CHT\_acc} & \textbf{CCL\_acc} & \textbf{EDT\_acc} \\ \hline
 Im001\_1 & 0 & 0 & 0 & 0 & 100.0 & 100.0 & 100.0 \\ 
 Im004\_1 & 4 & 2 & 3 & 2 & 50.0 & 75.0 & 50.0 \\ 
 Im015\_1 & 7 & 1 & 2 & 1 & 14.29 & 28.57 & 14.29 \\ 
 Im022\_1 & 15 & 8 & 10 & 8 & 53.33 & 66.67 & 53.33 \\ 
 Im037\_0 & 6 & 5 & 6 & 5 & 83.33 & 100.0 & 83.33 \\ 
 Im045\_0 & 39 & 24 & 25 & 24 & 61.54 & 64.1 & 61.54 \\ 
 Im050\_1 & 10 & 8 & 8 & 8 & 80.0 & 80.0 & 80.0 \\ 
 Im053\_1 & 12 & 9 & 12 & 9 & 75.0 & 100.0 & 75.0 \\ 
 Im069\_0 & 3 & 3 & 3 & 3 & 100.0 & 100.0 & 100.0 \\ 
 Im079\_0 & 0 & 2 & 3 & 2 & 0.0 & 0.0 & 0.0 \\ 
 Im095\_0 & 7 & 6 & 6 & 6 & 85.71 & 85.71 & 85.71 \\ 
 Im099\_0 & 37 & 26 & 26 & 26 & 70.27 & 70.27 & 70.27 \\ 
 Im108\_0 & 5 & 6 & 7 & 6 & 80.0 & 60.0 & 80.0 \\ \hline
 Total &   &   &   &   & 65.65 & 71.56 & 65.65 \\ 

\hline
\end{tabular}
\caption{Platelets counting results using 3 algorithms}
\label{Platelets counting results using 3 algorithms}
\end{table}


%TODO at the end we draw some tables that will contain the 3 

\section{Results}
\vspace{0.2in}
\hspace*{0.16in}

\begin{table}[H]
\centering
\begin{tabular}{|c|c|c|c|c|}
\hline
 \textbf{Image} & \textbf{Real\_Count} & \textbf{CCL} & \textbf{Watershed} & \textbf{CHT} \\ \hline
 \textbf{Im001\_0} & \textbf{21} & \textbf{23} & \textbf{19} & \textbf{17} \\ \hline
 \textbf{Im002\_1} & \textbf{11} & \textbf{13} & \textbf{19} & \textbf{10} \\ \hline
 \textbf{Im003\_0} & \textbf{20} & \textbf{17} & \textbf{13} & \textbf{11} \\ \hline

\end{tabular}
\caption{Results of each algorithm}
\label{Results of each algorithm}
\end{table}


